\documentclass[10pt,a4paper]{amsart}
\usepackage[margin=19mm]{geometry}
\usepackage{amsmath,amssymb,mathtools}
\usepackage[scr=rsfs,cal=euler]{mathalfa}
\usepackage{xfrac}
\usepackage[unicode,hidelinks,bookmarks=false]{hyperref}
\newcommand{\ie}{i.e.\ }
\newcommand{\cf}{cf.\ }
\newcommand{\resp}{respectively\ }
\newcommand{\Subsection}{\subsection}
\providecommand{\nobreakdash}{\nobreak\hbox{-}}
\setlength{\emergencystretch}{2em}
\multlinegap0pt
\usepackage{amssymb}
\usepackage[shortlabels]{enumitem}
\newtheorem{theorem}{Theorem}
\title{$G$-Exponential Families through Probability Coordinates}
\author{Manuela-Simona Cojocea}
\date{Source: arXiv:2609.03208v1 (CC BY 4.0)}
\begin{document}
\maketitle

\noindent\emph{Source and licence.} $G$-Exponential Families through Probability Coordinates. Version arXiv:2609.03208v1 (CC BY 4.0), distributed by arXiv under the Creative Commons Attribution 4.0 International licence, http://creativecommons.org/licenses/by/4.0/, as recorded in the arXiv licence metadata for this exact version and shown on the abstract page https://arxiv.org/abs/2609.03208v1. Full licence notice: \texttt{licenses/arxiv-2609.03208-CC-BY-4.0.txt}.\par\smallskip

\noindent\textit{Editorial scope.} Counted unit: the article's theorem thm:tail-equivalence (source0.tex lines 501--516). The article's complete original proof of it is delivered with it (source0.tex lines 518--536, one \texttt{proof} environment of 19 lines). No mathematics is written, repaired or re-derived: the statement and the proof are the article's own lines, in the article's own notation, and no retained line is substituted or re-flowed.  No further source-local context is needed: every cross-reference of every retained line resolves inside the two delivered spans.  Nothing was omitted from any retained span, so the package carries no omission record.  No retained line contains a citation, so the wrapper carries no bibliography and no bibliography entry.  No retained line contains a figure environment, an image-inclusion command or drawing code, so no image is delivered and the package declares no asset dependency.  Editorial formatting only, in addition: an AMS \texttt{amsart} preamble that restates the article's own declarations and macros used by the retained lines, namely the environments \texttt{theorem} on separate counters, together with the standard packages those lines need.  The retained environments print in this document as Theorem 1 in order of appearance, whereas the article prints the same items with the numbering of its own full document.  The complete native source of the article as distributed in the arXiv e-print for this exact version is bundled unchanged as \texttt{sources/209278/source0.tex}.
\begin{theorem}[Tail equivalence with the chart]\label{thm:tail-equivalence}
Let $f_\eta$ be a member of a $G$-exponential family, and suppose that the
probability-space density has a finite positive boundary limit,
\[
c_\eta:=\lim_{u\to1^-}p_\eta(u)\in(0,\infty),
\]
which holds in particular whenever $T$ extends continuously to $u=1$.
Denote by $F_\eta$ the distribution function of $f_\eta$. Then, as
$x\to\infty$,
\[
\frac{f_\eta(x)}{g(x)}\longrightarrow c_\eta
\qquad\text{and}\qquad
\frac{1-F_\eta(x)}{1-G(x)}\longrightarrow c_\eta .
\]
Every member of the family is therefore tail-equivalent to the chart law.
\end{theorem}

\begin{proof}
The first limit is immediate: $f_\eta(x)/g(x)=p_\eta(G(x))\to c_\eta$
since $G(x)\to1^-$. For the second, fix $\varepsilon\in(0,c_\eta)$ and
choose $\delta>0$ with $|p_\eta(u)-c_\eta|<\varepsilon$ for all
$u\in(1-\delta,1)$. For $x$ large enough that $G(x)>1-\delta$,
\[
1-F_\eta(x)=\int_x^\infty p_\eta(G(t))\,g(t)\,dt,
\]
and on the domain of integration $G(t)\ge G(x)>1-\delta$, so
\[
(c_\eta-\varepsilon)\int_x^\infty g(t)\,dt
\;\le\;
1-F_\eta(x)
\;\le\;
(c_\eta+\varepsilon)\int_x^\infty g(t)\,dt .
\]
Since $\int_x^\infty g=1-G(x)$, dividing by $1-G(x)$ and letting
$x\to\infty$, then $\varepsilon\to0$, gives the claim.
\end{proof}

\end{document}
